%

\section{1733 Giovanni Girolamo Saccheri} % lang-pl
\label{sekcja_saccheri}

Dopiero w 1733 roku publikacji doczekają się naukowe dociekania dotyczące piątego postulatu.
Girolamo Saccheri, lombardzki jezuita, tuż przed śmiercią wyda książkę ,,Euclides ab omni naevo vindicatus''\footnote{Euklides oczyszczony ze wszelkiej skazy.}.
Wprowadzi w niej:

\begin{definition}
[czworokąt Saccheriego] % lang-pl
	Czworokąt $ABCD$, w którym kąty $A$ i $B$ są proste, a boki $AD$ i $BC$ są równe. % lang-pl
	Odcinek $AB$ nazywamy podstawą, $DC$ szczytem. % lang-pl
\index{czworokąt!Saccheriego} % lang-pl
\end{definition}

Zafascynowany \emph{reductio ad absurdum}, zrobi co następuje.
Łatwo pokazać, że kąty przy wierzchołkach $C$, $D$ w czworokącie Saccheriego są równe.
Mogą być ostre, proste lub rozwarte.
Saccheri nazywa te możliwości kolejno hipotezą kąta ostrego, prostego i rozwartego i stara się pokazać, że jedynie ta środkowa nie prowadzi do sprzeczności.
Milcząco zakłada przy tym tak jak Euklides, że proste są nieskończone i szybko rozprawia się z kątem rozwartym.

Po żmudnym wyprowadzeniu licznych twierdzeń geometrii nieeuklidesowej (nieistniejącej wtedy jeszcze), Saccheri za wszelką cenę zechce dokończyć dowód nie wprost.
Wprowadzi na siłę niejasne elementy nieskończone zamiast przyznać, że nie potrafi znaleźć sprzeczności.

\color{red}
Wstąpił do zakonu w wieku 18 lat. W roku 1690 wyjechał na studia do Mediolanu, gdzie oprócz filozofii i teologii studiował również matematykę, zachęcony przez Tomasso Cevę, brata Giovanniego Cevy. W 1694 otrzymał święcenia kapłańskie.

Now considered an early exploration of non-Euclidean geometry, Euclides ab omni naevo vindicatus (Euclid Freed of Every Flaw) languished in obscurity until it was rediscovered by Eugenio Beltrami, in the mid-19th century.[9]

One tool that Saccheri developed in his work (now called a Saccheri quadrilateral) has a precedent in the 11th-century Persian polymath Omar Khayyám's Discussion of Difficulties in Euclid (Risâla fî sharh mâ ashkala min musâdarât Kitâb 'Uglîdis). Khayyam, however, made no significant use of the quadrilateral, whereas Saccheri explored its consequences deeply.[11]

% CLAUDE
Giovanni Girolamo Saccheri będzie uczniem Tommasa Cevy, brata Giovanniego, autora twierdzenia Cevy (w tej książce z numerem \ref{twierdzenie_cevy}). % lang-pl
W 1733 roku Giovanni Girolamo Saccheri spróbuje zbadać, co by było gdyby piąty postulat Euklidesa okazał się fałszywy. % lang-pl
tytuł przypomina słowa Savile'a o dwóch skazach (zobacz sekcję \ref{sekcja_piaty_postulat}) \cite[s. 305]{hartshorne_2000}. % lang-pl

Figura ta wróci do Claviusa, który w komentarzu do twierdzenia I.29 zaproponuje aksjomat równoodległości; wydanie Euklidesa Claviusa poleci Saccheriemu jego nauczyciel Tommaso Ceva, więc można przypuszczać, że to Clavius zainspiruje go do dalszych badań \cite[s. 306]{hartshorne_2000}. % lang-pl
Podobny czworokąt rozważą już wcześniej Sabit ibn Kurra (IX wiek), Omar Chajjam (XI wiek) i, tuż przed Saccherim, Giordano Vitale. % lang-pl

Wygląda na to, że książkę wkrótce po wydaniu wycofano ze sprzedaży, więc jej wpływ na współczesnych będzie znikomy \cite[s. 285]{eves1_1972}; Hartshorne dodaje, że praca Saccheriego, może zbyt wyprzedzająca swój czas, nie doczeka się uznania i przeleży w zapomnieniu do końca XIX wieku \cite[s. 306]{hartshorne_2000}. % lang-pl
Dziś wiadomo, że odrzucenie hipotezy kąta rozwartego było kruche: Saccheri (i po nim Legendre) wykluczy ją dzięki ciągłości, czyli aksjomatowi Archimedesa, a Max Dehn poda przykład nieArchimedesowej płaszczyzny Hilberta, w której ta hipoteza zachodzi \cite[s. 208--209]{greenberg_2010}. % lang-pl
Twierdzenie, że każda płaszczyzna Hilberta jest jednego z trzech typów (suma kątów trójkąta mniejsza, równa albo większa od $\pi$ we wszystkich trójkątach naraz), Saccheri miał już w 1733 roku, choć z użyciem ciągłości \cite[s. 208]{greenberg_2010}. % lang-pl

Dopiero w 1889 roku Eugenio Beltrami przypomni tę książkę światu (zobacz sekcję \ref{sekcja_beltrami_cayley_klein}). % lang-pl
% Źródła: Eves s. 285--286, 289, 304; Greenberg 2010, s. 208--209; Wikipedia: Johann Heinrich Lambert, Lambert quadrilateral.
% https://en.wikipedia.org/wiki/Lambert_quadrilateral


\color{black}

O Saccherim napisze Eves \cite[s. 284, 285, 298-301]{eves1_1972}.

%